\documentclass[aspectratio=169,11pt]{beamer}
% Shared CMPSC 480 Beamer preamble.
\usepackage[T1]{fontenc}
\usepackage[utf8]{inputenc}
\usepackage{lmodern}
\usepackage{amsmath,amssymb,mathtools}
\usepackage{booktabs,array,tabularx}
\usepackage{xcolor}
\usepackage{listings}
\usepackage{tikz}
\usetikzlibrary{arrows.meta,positioning,shapes.geometric}

% Ignore only negligible TeX box diagnostics that are below visual tolerance.
% Larger layout defects still appear in the build log and in rendered QA.
\vfuzz=3pt
\hbadness=2000

\definecolor{courseink}{HTML}{111111}
\definecolor{coursegray}{HTML}{EDEDED}
\definecolor{courserule}{HTML}{B8BCC4}
\definecolor{courseblue}{HTML}{3D8DFF}
\definecolor{courselightblue}{HTML}{D0EDFA}
\definecolor{coursemuted}{HTML}{5C6470}
\definecolor{coursegreen}{HTML}{0B7A53}
\definecolor{coursered}{HTML}{B3261E}

\usetheme{default}
\usefonttheme{professionalfonts}
\setbeamertemplate{navigation symbols}{}
\setbeamersize{text margin left=0.65cm,text margin right=0.65cm}

\setbeamercolor{normal text}{fg=courseink,bg=white}
\setbeamercolor{frametitle}{fg=courseink,bg=white}
\setbeamercolor{title}{fg=courseink,bg=white}
\setbeamercolor{subtitle}{fg=coursemuted,bg=white}
\setbeamercolor{structure}{fg=courseblue}
\setbeamercolor{alerted text}{fg=coursered}
\setbeamercolor{block title}{fg=courseink,bg=courselightblue}
\setbeamercolor{block body}{fg=courseink,bg=coursegray}
\setbeamercolor{example text}{fg=coursegreen}

\setbeamerfont{title}{size=\fontsize{25}{29}\selectfont,series=\bfseries}
\setbeamerfont{subtitle}{size=\fontsize{13}{16}\selectfont}
\setbeamerfont{frametitle}{size=\fontsize{19}{22}\selectfont,series=\bfseries}
\setbeamerfont{normal text}{size=\fontsize{11}{14}\selectfont}
\setbeamerfont{footline}{size=\fontsize{7.5}{9}\selectfont}

\setbeamertemplate{frametitle}{%
  \nointerlineskip
  % Use content-driven height so a long assessment title can wrap without
  % being clipped by a fixed-height title bar.
  \begin{beamercolorbox}[wd=\paperwidth,sep=0.17cm,leftskip=0.48cm,rightskip=0.48cm]{frametitle}
    \insertframetitle
  \end{beamercolorbox}
  \vspace{-0.08cm}
  {\color{courserule}\rule{\textwidth}{0.35pt}}
}

\setbeamertemplate{footline}{%
  \leavevmode
  \begin{beamercolorbox}[wd=\paperwidth,ht=2.6ex,dp=1.1ex,leftskip=.65cm,rightskip=.65cm]{author in head/foot}
    \color{coursemuted}\insertshorttitle\hfill\insertframenumber/\inserttotalframenumber
  \end{beamercolorbox}
}

\setbeamertemplate{itemize item}{\color{courseblue}\small$\blacktriangleright$}
\setbeamertemplate{itemize subitem}{\color{courseblue}\scriptsize$\bullet$}
\setbeamertemplate{enumerate item}{\color{courseblue}\insertenumlabel.}

\lstdefinestyle{coursepython}{
  language=Python,
  basicstyle=\ttfamily\fontsize{8.5}{10}\selectfont,
  keywordstyle=\color{courseblue}\bfseries,
  commentstyle=\color{coursemuted}\itshape,
  stringstyle=\color{coursegreen},
  showstringspaces=false,
  columns=fullflexible,
  keepspaces=true,
  frame=single,
  rulecolor=\color{courserule},
  backgroundcolor=\color{coursegray},
  breaklines=true,
  tabsize=4,
  xleftmargin=0.15cm,
  xrightmargin=0.15cm,
  framexleftmargin=0.15cm,
  framexrightmargin=0.15cm
}
\lstset{style=coursepython}

\newcommand{\points}[1]{\hfill{\color{courseblue}\textbf{[#1 points]}}}
\newcommand{\deliverable}[1]{\textbf{\color{courseblue}#1}}
\newcommand{\code}[1]{\texttt{#1}}
\newcommand{\keyidea}[1]{%
  \begin{beamercolorbox}[rounded=false,sep=0.55em,wd=\linewidth]{block body}
    \textbf{Key idea.} #1
  \end{beamercolorbox}
}
\newcommand{\answerlabel}{\textbf{\color{coursegreen}Answer.}\ }
\newcommand{\gradinglabel}{\textbf{\color{courseblue}Grading guidance.}\ }

\newenvironment{questionblock}[1]
  {\begin{block}{#1}}
  {\end{block}}

\newenvironment{solutionblock}[1]
  {\begin{block}{#1}}
  {\end{block}}

\AtBeginSection[]{
  \begin{frame}
    \vfill
    {\usebeamerfont{title}\color{courseink}\insertsectionhead\par}
    \vspace{0.4cm}
    {\color{courseblue}\rule{0.32\paperwidth}{3pt}}
    \vfill
  \end{frame}
}


% CMPSC480_TOTAL_POINTS: 10
% CMPSC480_ITEM_IDS: 1,2,3,4
\title[Week 5 Assignment Questions]{Week 5 Assignment: Tabular Q-Learning}
\subtitle{A tested model-free agent with controlled exploration and oracle comparison}
\author{CMPSC 480 - Student Question Deck}
\date{}

\begin{document}


\begin{frame}
  \titlepage
\end{frame}

\begin{frame}{Assessment overview}
\begin{columns}[T,onlytextwidth]
\begin{column}{0.62\textwidth}
\Large Train a tabular Q-learning agent and evaluate learning behavior against a planning oracle.

\vspace{0.35cm}
\normalsize
Coverage: Week 5: TD learning, exploration, terminal targets, and empirical evaluation
\end{column}
\begin{column}{0.33\textwidth}
\begin{block}{Points}10 total\end{block}
\begin{block}{Expected time}6-8 hours\end{block}
\begin{block}{Submission}Repository, learning curves, and report\end{block}
\end{column}
\end{columns}
\end{frame}

\begin{frame}{Learning objectives}
\begin{itemize}
  \item Separate environment dynamics from agent state.
  \item Implement epsilon-greedy action selection with deterministic evaluation.
  \item Apply the terminal-aware Q-learning update.
  \item Run reproducible training and evaluation phases.
  \item Compare a learned policy with value-iteration ground truth.
\end{itemize}
\end{frame}

\begin{frame}{Requirements checklist}
\small
\begin{itemize}
  \item Use Python 3.11 and NumPy only unless the prompt explicitly permits another package.
  \item Preserve the published interfaces and make all tie-breaking deterministic.
  \item State assumptions, validate inputs, and handle empty, terminal, and unreachable cases.
  \item Include focused tests whose names explain the defect each test would expose.
  \item Report measured evidence; do not replace experimental results with unsupported claims.
  \item Seed every source of randomness used in reported experiments.
\end{itemize}
\end{frame}

\begin{frame}{Task 1 - Agent and exploration contract (2 points)}
\small
Implement a Q-table agent without embedding environment transition logic.

\vspace{0.2cm}
\begin{enumerate}
\item[\textbf{a.}] Define observe, select\_action, update, and evaluate-mode behavior. \hfill{\color{courseblue}\textbf{[1 points]}}
\item[\textbf{b.}] Specify epsilon-greedy selection and a bounded decay schedule. \hfill{\color{courseblue}\textbf{[1 points]}}
\end{enumerate}
\end{frame}

\begin{frame}{Task 2 - Q-learning update (3 points)}
\small
Implement and validate the off-policy temporal-difference update.

\vspace{0.2cm}
\begin{enumerate}
\item[\textbf{a.}] State the update and define the TD error. \hfill{\color{courseblue}\textbf{[1 points]}}
\item[\textbf{b.}] Compute one update for Q=2, r=1, gamma=0.9, max next Q=4, and alpha=0.5. \hfill{\color{courseblue}\textbf{[1 points]}}
\item[\textbf{c.}] Explain and test terminal handling. \hfill{\color{courseblue}\textbf{[1 points]}}
\end{enumerate}
\end{frame}

\begin{frame}{Task 3 - Training protocol (2 points)}
\small
Build a reproducible episodic training loop with safety limits.

\vspace{0.2cm}
\begin{enumerate}
\item[\textbf{a.}] Describe reset, interaction, update, and episode termination order. \hfill{\color{courseblue}\textbf{[1 points]}}
\item[\textbf{b.}] Log return, episode length, epsilon, and moving-average return. \hfill{\color{courseblue}\textbf{[1 points]}}
\end{enumerate}
\end{frame}

\begin{frame}{Task 4 - Evaluation and tests (3 points)}
\small
Measure policy quality and diagnose sensitivity to exploration and rewards.

\vspace{0.2cm}
\begin{enumerate}
\item[\textbf{a.}] Evaluate the greedy policy without learning or exploration. \hfill{\color{courseblue}\textbf{[1 points]}}
\item[\textbf{b.}] Compare greedy actions or returns with a value-iteration oracle. \hfill{\color{courseblue}\textbf{[1 points]}}
\item[\textbf{c.}] Run one controlled hyperparameter ablation and cover sparse-reward and unseen-state cases. \hfill{\color{courseblue}\textbf{[1 points]}}
\end{enumerate}
\end{frame}

\begin{frame}{Edge cases and defensive programming}
\small
\begin{itemize}
  \item No legal actions are returned.
  \item The transition is terminal.
  \item All action values tie.
  \item A state has never been visited.
  \item An episode reaches the safety step limit.
\end{itemize}
\end{frame}

\begin{frame}{Submission instructions}
\small
\begin{itemize}
  \item Submit source, tests, seeds, and hyperparameters.
  \item Include learning curves plus a compact oracle-comparison table.
  \item Document the exact command that reproduces the reported run.
\end{itemize}
\vspace{0.2cm}
\alert{Do not submit credentials, generated caches, or unapproved libraries.}
\end{frame}

\begin{frame}{Grading rubric}
\scriptsize
\begin{tabularx}{\textwidth}{>{\bfseries}p{0.28\textwidth} p{0.08\textwidth} X}
\toprule
Category & Points & Required evidence \\
\midrule
Agent and exploration contract & 2 & Agent API, epsilon schedule, and action-selection tests. \\
Q-learning update & 3 & Update function and hand-computed unit tests. \\
Training protocol & 2 & Training loop, seeds, and logged metrics. \\
Evaluation and tests & 3 & Oracle comparison, ablation, and edge-case tests. \\
\bottomrule
\end{tabularx}
\end{frame}

\begin{frame}{Reflection prompt}
In 150-200 words:

Explain one gap between a learned policy and the planning oracle. Distinguish modeling mismatch from sampling or optimization error.
\end{frame}

\end{document}
